Monte Carlo (MC) simulation is used to estimate the signal acceptance
and to model the contributions from Standard Model background processes.
All simulated samples correspond to the ATLAS Run~2 simulation campaign
(MC16a, MC16d and MC16e), reproducing the detector and pileup conditions
of the data-taking periods.

After the event selection, the dominant background arises from
$t\bar{t}$ production in association with additional jets.
Subleading contributions originate from single-top production,
$W$+jets and $Z$+jets processes, diboson production ($WW$, $WZ$, $ZZ$),
and associated production of a top-quark pair with a vector or Higgs boson
($t\bar{t}V$ and $t\bar{t}H$).
Multijet production can contribute through the misidentification of jets
as leptons or through non-prompt leptons originating from heavy-flavour
hadron decays, but its contribution is found to be negligible after the
event selection.

Both signal and background samples include the simulation of pileup
interactions and are processed through either the full \textsc{Geant4}
detector simulation or the ATLAS fast simulation framework (AFII).
All simulated events are reconstructed with the same software chain as
used for the data.

Further details about the modelling of each background process are
summarised below.

\begin{description}

\item[$t\bar{t}$ + jets:] 

The nominal $t\bar{t}$ sample is generated using the
\textsc{Powheg-Box}~v2 generator at next-to-leading order (NLO)
in QCD with the \textsc{NNPDF3.0} parton distribution function (PDF)
set in the five-flavour scheme (5FS).
The parameter $h_{\mathrm{damp}}$, which regulates the transverse
momentum of the first emission beyond the Born configuration,
is set to $1.5\,m_t$ with $m_t = 172.5$~GeV.
Parton showering and hadronisation are simulated using
\textsc{Pythia}~8 with the A14 tune.

Additional filtered samples are used to enhance the statistical
precision in regions with large jet multiplicity or high
$b$-tag multiplicity.

An additional $t\bar{t}+\geq 1b$ sample is generated using
\textsc{Powheg-Box-Res} interfaced with \textsc{OpenLoops}
in the four-flavour scheme (4FS).
This sample is used to evaluate modelling uncertainties
associated with heavy-flavour production in $t\bar{t}$ events.

Alternative samples are used to evaluate modelling uncertainties
related to the hard-scattering generator, parton shower and QCD
radiation.
These include samples generated with
\textsc{MadGraph5\_aMC@NLO} interfaced with \textsc{Pythia8}
to assess the impact of the matrix-element generator,
and samples where the \textsc{Powheg} matrix element is interfaced
to \textsc{Herwig7} to evaluate the dependence on the
parton-shower and hadronisation model.
Additional variations of the radiation parameters in the
A14 tune are also used to probe uncertainties related to
initial- and final-state radiation.

The $t\bar{t}$ events are categorised according to the flavour
content of additional jets not originating from the top-quark decays,
separating heavy-flavour ($t\bar{t}+b\bar{b}$, $t\bar{t}+c\bar{c}$)
and light-flavour components.

The samples are normalised to the theoretical cross section
computed with \textsc{Top++} at NNLO in QCD including NNLL
soft-gluon resummation, corresponding to
$\sigma_{t\bar{t}} = 832^{+46}_{-51}$~pb.

\item[$t\bar{t}t\bar{t}$:]

The production of four top quarks is simulated using
\textsc{MadGraph5\_aMC@NLO} at NLO accuracy in QCD with the
NNPDF3.1 NLO PDF set.
The renormalisation and factorisation scales are chosen as
$\mu_r=\mu_f=m_T/4$, where $m_T$ is defined as the scalar sum of
the transverse masses of the particles produced in the
matrix-element calculation.

Top-quark decays are handled using \textsc{MadSpin}
to preserve spin correlations.
Parton showering and hadronisation are simulated with
\textsc{Pythia}~8 using the A14 tune.
Heavy-flavour hadron decays are modelled using
\textsc{EvtGen}.

An additional sample generated with \textsc{Herwig7}
is used to evaluate uncertainties associated with the
parton shower and hadronisation model.
Variations of the renormalisation and factorisation scales
in the matrix-element calculation are used to estimate
uncertainties related to missing higher-order corrections.

All $t\bar{t}t\bar{t}$ samples are normalised to the
NLO(QCD+EW)+NLL$'$ theoretical prediction of
$\sigma_{t\bar{t}t\bar{t}} = 13.37$~fb.

\item[$t\bar{t}V$:]

The associated production of a top-quark pair with a vector boson
($t\bar{t}W$, $t\bar{t}Z$) is simulated using
\textsc{MadGraph5\_aMC@NLO} with matrix elements including
up to two additional partons at leading order.
The events are interfaced with \textsc{Pythia}~8 for
parton showering and hadronisation using the A14 tune.
Alternative samples generated with \textsc{Sherpa}
are used to evaluate modelling uncertainties associated
with the matrix-element generator.

\item[$t\bar{t}H$:]

The production of $t\bar{t}H$ events is simulated using
the \textsc{Powheg}~v2 generator at NLO accuracy in QCD
with the NNPDF3.0 NLO PDF set.
The $h_{\mathrm{damp}}$ parameter is set to
$1.5\times (2m_t+m_H)/2$ with $m_H=125$~GeV.
Parton showering and hadronisation are modelled with
\textsc{Pythia}~8 using the A14 tune.
The sample is normalised to the theoretical cross section
computed at NLO in QCD including leading electroweak corrections.

An alternative sample generated with
\textsc{MadGraph5\_aMC@NLO} and interfaced with
\textsc{Pythia8} is used to evaluate uncertainties
associated with the choice of matrix-element generator.

\item[Single-top:]

Single-top production in the $tW$, $t$-channel and $s$-channel
processes is simulated using the \textsc{Powheg}~v2 generator
at NLO in QCD with the NNPDF3.0 NLO PDF set.
The $tW$ process is generated using the diagram removal scheme
to avoid overlap with $t\bar{t}$ production.
Parton showering and hadronisation are simulated with
\textsc{Pythia}~8 using the A14 tune.

Alternative samples generated with
\textsc{MadGraph5\_aMC@NLO} and samples interfaced
with \textsc{Herwig7} are used to estimate modelling
uncertainties associated with the generator and
parton-shower description.

\item[$V$+jets:]

The production of $W$+jets and $Z$+jets events is simulated
using the \textsc{Sherpa}~v2.2 generator.
Matrix elements are calculated at NLO accuracy for up to two
additional partons and at LO accuracy for up to four additional
partons.
They are matched and merged with the \textsc{Sherpa} parton
shower using the MEPS@NLO prescription.
The NNPDF3.0 NNLO PDF set is used.

\item[$VV$:]

Diboson production ($WW$, $WZ$, $ZZ$) is simulated using
\textsc{Sherpa}~v2.2.
Matrix elements are calculated at NLO accuracy in QCD
for up to one additional parton and at LO accuracy for up to
three additional partons.
The samples are matched and merged with the
\textsc{Sherpa} parton shower using the MEPS@NLO prescription
and the NNPDF3.0 NNLO PDF set.

\end{description}

The Monte Carlo generators used to simulate the signal and background
processes are summarised in Table~\ref{tab:MC_generators}.
Alternative samples used to evaluate modelling uncertainties
are also indicated.

\begin{table}[htbp]
\centering
\footnotesize
\setlength{\tabcolsep}{3pt}
\renewcommand{\arraystretch}{1.15}

\begin{tabular}{
p{2.0cm}
p{4.6cm}
p{1.9cm}
p{2.4cm}
p{1.9cm}
p{1.2cm}
}

\toprule
Process & Generator & ME order & PDF & Parton shower & Tune \\
\midrule

\multicolumn{6}{c}{\textbf{Signal processes}} \\
\midrule

$t\bar t H/A \rightarrow t\bar t t\bar t$
& \textsc{MadGraph5\_aMC@NLO}
& LO
& NNPDF3.1LO
& Pythia8
& A14 \\

& (\textsc{MadGraph5\_aMC@NLO})
& (LO)
& (NNPDF3.1LO)
& (Herwig7)
& (H7UE) \\

$S_8S_8 \rightarrow t\bar t t\bar t$
& \textsc{MadGraph5\_aMC@NLO} + MadSpin
& LO
& NNPDF3.1LO
& Pythia8
& A14 \\

& (\textsc{MadGraph5\_aMC@NLO} + MadSpin)
& (LO)
& (NNPDF3.1LO)
& (Herwig7)
& (H7UE) \\

\midrule
\multicolumn{6}{c}{\textbf{Background processes}} \\
\midrule

\multirow{4}{*}{$t\bar t$ + jets}

& PowhegBox
& NLO
& NNPDF3.0NLO
& Pythia8
& A14 \\

& (\textsc{MadGraph5\_aMC@NLO})
& (NLO)
& (NNPDF3.0NLO)
& (Pythia8)
& (A14) \\

& (PowhegBox)
& (NLO)
& (NNPDF3.0NLO)
& (Herwig7)
& (H7UE) \\

& (PowhegBox, rad.\ variations)
& (NLO)
& (NNPDF3.0NLO)
& (Pythia8)
& (A14 Var3c) \\

$t\bar t +\geq 1b$ (4FS)
& PowhegBox-Res + OpenLoops
& NLO
& NNPDF3.0nlo nf4
& Pythia8
& A14 \\

\multirow{2}{*}{$t\bar t t\bar t$}

& \textsc{MadGraph5\_aMC@NLO}
& NLO
& NNPDF3.1NLO
& Pythia8
& A14 \\

& (\textsc{MadGraph5\_aMC@NLO})
& (NLO)
& (NNPDF3.1NLO)
& (Herwig7)
& (H7UE) \\

\multirow{2}{*}{$t\bar tV$}

& \textsc{MadGraph5\_aMC@NLO}
& LO
& NNPDF2.3LO
& Pythia8
& A14 \\

& (Sherpa 2.2.1)
& (MEPS@NLO)
& (NNPDF3.0NNLO)
& (Sherpa)
& (Sherpa) \\

\multirow{2}{*}{$t\bar tH$}

& PowhegBox
& NLO
& NNPDF3.0NLO
& Pythia8
& A14 \\

& (\textsc{MadGraph5\_aMC@NLO})
& (NLO)
& (NNPDF3.0NLO)
& (Pythia8)
& (A14) \\

\multirow{3}{*}{Single-top}

& PowhegBox ($tW$)
& NLO
& NNPDF3.0NLO
& Pythia8
& A14 \\

& (PowhegBox)
& (NLO)
& (NNPDF3.0NLO)
& (Herwig7)
& (H7UE) \\

& (\textsc{MadGraph5\_aMC@NLO})
& (NLO)
& (NNPDF3.0NLO)
& (Pythia8)
& (A14) \\

$V+$jets ($W/Z$)
& Sherpa 2.2.1
& MEPS@NLO
& NNPDF3.0NNLO
& Sherpa
& Sherpa \\

$VV$ ($WW$, $WZ$, $ZZ$)
& Sherpa 2.2.1/2.2.2
& MEPS@NLO
& NNPDF3.0NNLO
& Sherpa
& Sherpa \\

\bottomrule
\end{tabular}

\caption{Monte Carlo generators used to simulate the signal and
background processes considered in this analysis.
Alternative samples used to evaluate modelling uncertainties
are shown in parentheses.}

\label{tab:MC_generators}
\end{table}